\section{TEE Design Space and Evidence Boundaries}

Before comparing Arm CCA, RISC-V CoVE/AP-TEE, and confidential I/O, it is
useful to separate the design dimensions that recur across TEE systems. Survey
and SoK sources provide the taxonomy for this section, while platform-specific
mechanism claims in later sections still return to primary papers and official
specifications~\cite{schneider2022soktee,li2024sokteechoices,boubakri2025riscvtee,sok-tee}.

\begin{table*}[t]
\centering
\caption{TEE Design-Space Dimensions and Evidence Boundaries}
\label{tab:tee_design_space}
\begin{tabular}{@{}p{0.15\linewidth}p{0.29\linewidth}p{0.24\linewidth}p{0.22\linewidth}@{}}
\toprule
\textbf{Dimension} & \textbf{Survey question} & \textbf{Typical mechanisms} & \textbf{Evidence boundary} \\ \midrule
Launch and admission & How does a protected domain enter a known initial state before receiving secrets? & Secure boot, measured launch, Realm/TVM creation, monitor-mediated acceptance & Use SoK sources for vocabulary; use CCA, AP-TEE, and platform specs for mechanism details. \\
Runtime isolation & Which agents can observe or mutate execution state and protected memory during operation? & World/domain separation, PMP, GPT/GPC, page lifecycle, enclave monitor policy & Distinguish process, enclave, Realm, and confidential-VM threat models. \\
Trusted I/O & How can a protected workload use devices without handing data authority back to the host? & SMMU/IOMMU/IOPMP, SPDM, TDISP, PCIe IDE, CoVE-IO, device assignment & Treat drafts, gated specifications, and vendor documents under their restricted evidence classes. \\
Secure storage and persistence & What protects data once it leaves CPU-local protected memory or survives reboot? & Sealing, RPMB/fTPM, encrypted storage, freshness metadata, remote-storage integrity & Do not infer storage freshness or rollback protection from execution isolation alone. \\
Attestation and verifier policy & What evidence lets a remote party decide whether to release keys or trust a service? & RATS roles, EAT tokens, platform evidence, device evidence, endpoint binding & RFCs define roles and token formats; platform claims need platform-specific evidence. \\
TCB and excluded attacks & Which hardware, firmware, runtime, host, and device components remain trusted or out of scope? & Monitor/RMM/TSM, firmware roots, device managers, verifier policy, side-channel exclusions & Keep side-channel and physical attacks as threat-boundary statements unless a mechanism directly addresses them. \\ \bottomrule
\end{tabular}
\end{table*}

Table~\ref{tab:tee_design_space} provides the bridge used by the rest of the
survey. Arm CCA and RISC-V CoVE/AP-TEE both address launch, runtime isolation,
attestation, and TCB control, but they expose different interfaces and
standardization status. Confidential I/O and secure storage extend the same
questions beyond CPU execution into DMA, interrupts, links, devices, and
persistent or remote state. Keeping these dimensions separate prevents three
common overclaims: treating attestation as isolation, treating memory ownership
as storage freshness, and treating device identity as complete trusted I/O.

Deployment substrates introduce another boundary that is easy to overstate.
\textbf{CoFunc} studies confidential serverless functions through split
containers, placing a function-oriented microkernel and library OS inside a CVM
while leaving Linux container management in an untrusted shadow container
outside the CVM~\cite{shi2025cofunc}. This is useful evidence for confidential
container and serverless packaging because it exposes the tension between CVM
cold start, resource sharing, TCB size, and function isolation. It is not an
Arm CCA mechanism paper and should not be used to justify Realm lifecycle,
RMM/RMI semantics, device assignment, or trusted-I/O claims.
